\section{Consequentialism}

This content comes from Chapter 9 of \textit{The Fundamentals of Ethics}.

\subsection{Nature of Consequentialism}

\begin{readingnotes}
	\textbf{Consequentialism} is a family of theories united by the belief that acts are morally right just because they produce the best overall results.

	The most prominent version is \textbf{Act Utilitarianism}, which states that well-being is the only thing that is intrinsically valuable, and an action is morally required just because it does more to improve overall well-being than any other action you could have done in the circumstances.
\end{readingnotes}

\begin{keyquote}[The Consequentialist Motto]
	Do all the good you can, by all the means you can, in all the ways you can, in all the places you can, at all the times you can, to all the people you can, as long as ever you can.
\end{keyquote}

\begin{idea}
	\textbf{Optimific}: An action is optimific if it yields the greatest balance of benefits over drawbacks.

	\textbf{Principle of Utility}: The ultimate moral standard that requires us to maximize overall well-being.
\end{idea}

\subsection{Determining the Right Action}

\begin{readingnotes}
	Utilitarianism provides a five-step process to determine if an act is optimific:
	\begin{enumerate}
		\item Identify what is \textbf{intrinsically good} (valuable for its own sake, e.g., happiness).
		\item Identify what is \textbf{intrinsically bad} (e.g., pain, misery).
		\item Determine all available \textbf{options} at the moment.
		\item For each option, determine the value of its results (good vs. bad).
		\item Pick the action that yields the \textbf{greatest net balance} of good over bad.
	\end{enumerate}

	\textbf{Misunderstandings:}
	\begin{itemize}
		\item We do \textit{not} always have to benefit the greatest number of people (sometimes a huge benefit to a few outweighs a small benefit to many).
		\item We do \textit{not} always choose the act with the most happiness (we must minimize misery as well; net balance is what counts).
	\end{itemize}
\end{readingnotes}

\subsection{Actual vs. Expected Results}

\begin{readingnotes}
	There is a debate regarding moral knowledge and the future:
	\begin{itemize}
		\item \textbf{Standard View (Actual Results):} Actions are right only if they \textit{actually} bring about the best results. If you intend good but cause harm, the action was immoral, though you may not be blameworthy.
		\item \textbf{Alternative View (Expected Results):} Acts are morally right if they are \textit{reasonably expected} to be optimific.
	\end{itemize}

	Most utilitarians separate the evaluation of the \textbf{action} (based on results) from the \textbf{intention} (based on expectation). You can be praised for good intentions even if the action was wrong (failed to be optimific).
\end{readingnotes}

\subsection{Attractions of Utilitarianism}

\begin{readingnotes}
	\begin{itemize}
		\item \textbf{Impartiality:} Everyone's well-being counts equally, regardless of status, wealth, or race.
		\item \textbf{Conflict Resolution:} It offers a single ultimate rule (maximize well-being) to resolve complex moral dilemmas.
		\item \textbf{Moral Flexibility:} No moral rule is absolute (except the principle of utility). Traditional rules (don't lie, don't kill) can be broken if doing so maximizes well-being.
		\item \textbf{The Moral Community:} Utilitarianism expands the moral community to include any being that can \textbf{suffer}, including non-human animals.
	\end{itemize}
\end{readingnotes}

\begin{argument}
	\textbf{The Argument from Marginal Cases}
	\begin{enumerate}
		\item If it is immoral to kill and eat "marginal" human beings, and to painfully experiment on them, then it is immoral to treat nonhuman animals this way.
		\item It is (almost) always immoral to kill and eat "marginal" human beings, and to painfully experiment on them.
		\item Therefore, it is (almost) always immoral to kill and eat animals, and to painfully experiment on them.
	\end{enumerate}
	\textit{(Note: "Marginal" humans refers to those with mental lives no more developed than animals, e.g., severe brain trauma or infants. The argument claims that since they can suffer, they are moral equals to animals.)}
\end{argument}

\subsection{Slippery Slope Arguments}

\begin{readingnotes}
	Utilitarianism is often used to critique social change via \textbf{Slippery Slope Arguments}.
	These arguments oppose a policy (like voluntary euthanasia) by predicting that even if the immediate results are good, it will inevitably lead to a future disaster (e.g., killing the unwilling).

	To defeat a slippery slope argument, one must show that the "dire forecast" at the bottom of the slope is unreasonable or unlikely.
\end{readingnotes}
